% (c) Nikita Lisitsa, lisyarus@gmail.com, 2025

\documentclass[10pt]{beamer}

\usepackage[T2A]{fontenc}
\usepackage[russian]{babel}
\usepackage{minted}

\usepackage[most]{tcolorbox}
\tcbuselibrary{minted}

\usepackage{graphicx}
\graphicspath{ {./images/} }

\usetheme{metropolis}

\usepackage{adjustbox}

\usepackage{color}
\usepackage{soul}

\usepackage{hyperref}

\definecolor{codebg}{RGB}{29,35,49}

\setminted{fontsize=\footnotesize}

\setlength\partopsep{-\topsep}

\definecolor{blue}{rgb}{0,0,1}
\definecolor{red}{rgb}{1,0,0}

\usemintedstyle{tango}

\makeatletter
\newcommand{\slideimage}[1]{
  \begin{figure}
    \begin{adjustbox}{width=\textwidth, totalheight=\textheight-2\baselineskip-2\baselineskip,keepaspectratio}
      \includegraphics{#1}
    \end{adjustbox}
  \end{figure}
}
\makeatother

\title{Компьютерная графика}
\subtitle{Практика 5: Bind groups, uniform буферы, текстуры, сэмплеры}
\date{2026}

\setbeamertemplate{footline}[frame number]

\begin{document}

\frame{\titlepage}

\begin{frame}[fragile]
\frametitle{Практика 5}
\begin{itemize}
\item Начиная с этой практики можно вовсю пользоваться математическими библиотеками (\mintinline{rust}|glam| в Rust, \mintinline{cpp}|math| в C++)
\item \alert{\textbf{NB}}: в Rust glam матрицы хранятся по столбцам (как и ожидает WebGPU), в C++ math -- по строкам (надо их транспонировать, либо умножать в шейдере с другой стороны)
\end{itemize}
\end{frame}

\begin{frame}[fragile]
\frametitle{Практика 5}
\slideimage{0.png}
\end{frame}

\begin{frame}[fragile]
\frametitle{Задание 1}
Создаём uniform buffer и bind группу
\begin{itemize}
\item Размер буфера -- как у immediates, \mintinline{cpp}|usage = Uniform|
\item Создаём bind group layout с одним entry: \mintinline{cpp}|binding = 0|, \mintinline{cpp}|visibility = Vertex|, \mintinline{cpp}|buffer.type = Uniform|
\item Создаём bind group с этим layout, тоже с одним entry: \mintinline{cpp}|binding = 0|, \mintinline{cpp}|buffer =| наш буфер
\item \textbf{\alert{NB}}: Пока ничего не изменится, мы никак не используем этот буфер :)
\end{itemize}
\end{frame}

\begin{frame}[fragile]
\frametitle{Задание 2}
Используем uniform buffer
\begin{itemize}
\item Убираем \mintinline{cpp}|immediateSize| из pipeline layout'а
\item Добавляем наш bind group layout в pipeline layout
\item Запись immediates в render pass'е заменяем на запись в наш буфер (\mintinline{cpp}|queue.write_buffer|) \textbf{перед созданием} render pass'а (в том же формате: сначала матрица \mintinline{cpp}|model|, затем \mintinline{cpp}|viewProjection|)
\item Добавляем \mintinline{cpp}|render_pass.set_bind_group(0, bind_group)| перед вызовом \mintinline{cpp}|draw_indexed()|
\item В шейдере заменяем \mintinline{wgsl}|var<immediate>| на \mintinline{wgsl}|@group(0) @binding(0) var<uniform>|
\item Если всё сделать правильно, картинка не изменится -- мы берём те же матрицы, но теперь из другого места
\end{itemize}
\end{frame}

\begin{frame}[fragile]
\frametitle{Задание 2}
\slideimage{0.png}
\end{frame}

\begin{frame}[fragile]
\frametitle{Задание 3}
Добавляем текстурные координаты
\begin{itemize}
\item Добавляем в render pipeline ещё один атрибут вершины, \mintinline{cpp}|shaderLocation = 2|, формат \mintinline{cpp}|vec2f| (в массиве вершин он уже есть)
\item Передаём его из вершинного шейдера во фрагментный
\item Выводим его в качестве цвета, например как \mintinline{wgsl}|vec4f(texcoord, 0.0, 1.0)|
\end{itemize}
\end{frame}

\begin{frame}[fragile]
\frametitle{Задание 3}
\slideimage{3.png}
\end{frame}

\begin{frame}[fragile]
\frametitle{Задание 4}
Загружаем текстуру
\begin{itemize}
\item Создаём текстуру: \mintinline{cpp}|device.create_texture()|, \mintinline{cpp}|dimension = 2D|, \mintinline{cpp}|format = RGBA8Unorm|, \mintinline{cpp}{usage = CopyDst | TextureBinding}, размер взять из \mintinline{cpp}|cow_image|
\item Загружаем пиксели в нашу текстуру: \mintinline{cpp}|queue.write_texture()|, \mintinline{cpp}|aspect = All|, \mintinline{cpp}|bytes_per_row = cow_image.width * 4|, \mintinline{cpp}|rows_per_image = cow_image.height|
\item Создаём texture view (дескриптор можно оставить пустым -- все дефолтные настройки нам подойдут)
\item \textbf{\alert{NB}}: Картинка не изменится, мы никак не используем текстуру :)
\end{itemize}
\end{frame}

\begin{frame}[fragile]
\frametitle{Задание 5}
Создаём сэмплер и используем текстуру
\begin{itemize}
\item Добавляем ещё два entry в bind group layout: 
\begin{itemize}
\item \mintinline{cpp}|binding = 1|, \mintinline{cpp}|visibility = Fragment|, \mintinline{cpp}|texture.view_dimension = 2D|, \mintinline{cpp}|texture.sample_type = Float|
\item \mintinline{cpp}|binding = 2|, \mintinline{cpp}|visibility = Fragment|, \mintinline{cpp}|sampler.type = Filtering|
\end{itemize}
\item Создаём sampler: \mintinline{cpp}|device.create_sampler()|, \mintinline{cpp}|address_mode = Repeat|, \mintinline{cpp}|mag_filter = Nearest|, \mintinline{cpp}|min_filter = linear|, \mintinline{cpp}|mipmap_filter = Nearest|
\item Добавляем texture view и sampler в bind группу
\item В шейдере:
\begin{itemize}
\item Возвращаем нормальный цвет вместо текстурных координат
\item Добавляем текстуру и сэмплер
\item Записываем результат вызова \mintinline{wgsl}|textureSample(...)| в переменную \mintinline{wgsl}|albedo|
\end{itemize}
\end{itemize}
\end{frame}

\begin{frame}[fragile]
\frametitle{Задание 5}
\slideimage{5.png}
\end{frame}

\begin{frame}[fragile]
\frametitle{Задание 6}
Загружаем mipmaps
\begin{itemize}
\item При создании текстуры указываем \mintinline{cpp}|mip_level_count = 4|
\item Нулевой уровень -- сама основная текстура, она уже загружена
\item Заполним оставшиеся 3 уровня красным, зелёным, и синим цветами, соответственно
\item Каждый следующий уровень должен быть в 2 раза меньше предыдущего по обеим осям!
\item Подвигайте камеру, чтобы увидеть, как включаются разные mipmap-уровни
\end{itemize}
\end{frame}

\begin{frame}[fragile]
\frametitle{Задание 6}
\slideimage{6.png}
\end{frame}

\begin{frame}[fragile]
\frametitle{Задание 7*}
Анимируем текстуру
\begin{itemize}
\item Добавим что-нибудь к текстурным координатам во фрагментном шейдере, чтобы текстура двигалась со временем
\item Нужно будет как-то передать время в шейдер: можно как часть uniform-buffer'а, можно как immediate
\end{itemize}
\end{frame}

\begin{frame}[fragile]
\frametitle{Задание 7*}
\slideimage{7.png}
\end{frame}

\end{document}
